\partkey{RnとC}

\phantomsection
\part{\texorpdfstring{$\mathbb{R}^n$と$\mathbb{C}$}{Rn と C}}

\begin{abstract}
この章では，$n$次元数空間 $\mathbb{R}^n$ と複素数体 $\mathbb{C}$ を導入する．まず直積集合として $\mathbb{R}^n$ を定義し，ベクトルの和・スカラー倍・内積といった演算とその性質を確認する．次に，$\mathbb{R}^2$ に加法と乗法を定めることで複素数体 $\mathbb{C}$ を構成し，虚数単位 $i$ が $i^2 = -1$ を満たすことを示す．最後に，$\mathbb{R}^n$ の点列を定義し，点列の収束・コーシー列といった概念が成分ごとの議論に帰着されることを確認する．
\end{abstract}

\draftpart

\phantomsection
\section{数空間と複素数体}

\phantomsection
\subsection{\texorpdfstring{$n$次元数空間}{n次元数空間}}

\begin{definition}{$n$次元数空間}{n次元数空間}
  直積集合
  \[
    \mathbb{R} \times \mathbb{R} \times \dotsm \times \mathbb{R} = \mathbb{R}^n
  \]
  を$n$次元数空間といい，その元を$n$次元数ベクトルという．
\end{definition}

新しい対象を定義したら，次に「2つの対象がいつ等しいとみなされるのか」を約束しておく必要がある．ベクトルについては，次のように定める．

\begin{definition}{ベクトルの相等}{ベクトルの相等}
  2つのベクトル
  \[
    \bm{x} = (x_1,x_2,\ldots,x_n) \in \mathbb{R}^n,\quad \bm{y} = (y_1,y_2,\ldots,y_n) \in \mathbb{R}^n
  \]
  が等しいことを，任意の$i \in \{ 1, 2 , \ldots ,n \}$に対して$x_i = y _i$であることとする．
\end{definition}

\begin{definition}{ベクトルの和}{ベクトルの和}
  2つのベクトル
  \[
    \bm{x} = (x_1,x_2,\ldots,x_n) \in \mathbb{R}^n,\quad \bm{y} = (y_1,y_2,\ldots,y_n) \in \mathbb{R}^n
  \]
  の和を，
  \[
    \bm{x}+\bm{y} \coloneqq (x_1+y_1,x_2+y_2,\ldots,x_n+y_n) \in \mathbb{R}^n
  \]
  とする．
\end{definition}

\begin{remark}[記号$+$の使い分け]
  左辺の$+$は演算$+ \colon \mathbb{R}^n \times \mathbb{R}^n \to \mathbb{R}^n$を表しているが，
  右辺の$+$は演算$+ \colon  \mathbb{R} \times \mathbb{R} \to \mathbb{R}$を表しており，これを区別することなく演算記号を約束した．
\end{remark}

\begin{definition}{ベクトルのスカラー倍}{ベクトルのスカラー倍}
  ベクトル
  \[
    \bm{x} = (x_1,x_2,\ldots,x_n) \in \mathbb{R}^n
  \]
  と実数$ a \in \mathbb{R}$の積を，
  \[
    a  \cdot  \bm{x} \coloneqq (a \cdot  x_1,a\cdot  x_2,\ldots,a \cdot  x_n) \in \mathbb{R}^n
  \]
  とする．
\end{definition}

\begin{remark}[記号$\cdot$の使い分けと省略]
  ここでも，$a \cdot \bm{x}$の$ \cdot $は演算$\cdot  \colon \mathbb{R} \times   \mathbb{R}^n \to \mathbb{R}^n$を表しているが，
  右辺の$\cdot $は演算$\cdot \colon \mathbb{R} \times  \mathbb{R} \to \mathbb{R}$を表しており，これを区別することなく演算記号を約束した．
  また，次からは$\cdot $を省略することがある．
\end{remark}

\newpage

\begin{theorem}{ベクトルの演算に関する性質}{ベクトルの演算に関する性質}
  $\mathbb{R}^n$の任意のベクトル$\bm{x}$，$\bm{y}$，$\bm{z}$と$a \in \mathbb{R}$に対して，次の性質が成り立つ．
  \begin{enumerate}
    \item $\bm{x}+\bm{y} = \bm{y}+\bm{x}$
    \item $(\bm{x}+\bm{y})+\bm{z} = \bm{x}+(\bm{y}+\bm{z})$
    \item $\bm{x}+\bm{0} = \bm{x}$
    \item $\bm{x}+(-\bm{x}) = \bm{0}$
    \item $1 \cdot \bm{x} = \bm{x}$
    \item $a \cdot (b \cdot \bm{x}) = (ab) \cdot \bm{x}$
    \item $(a+b) \cdot \bm{x} = a \cdot \bm{x} + b \cdot \bm{x}$
    \item $a \cdot (\bm{x}+\bm{y}) = a \cdot \bm{x} + a \cdot \bm{y}$
  \end{enumerate}
\end{theorem}

\begin{leftbar}
  \begin{proof}
    1，2，5，6は演算の定義から明らか．
    3，4は，$\bm{0} = (0,0,\ldots,0)$，$-\bm{x} = (-x_1,-x_2,\ldots,-x_n)$と定義すればよい．
    7，8は，$a \cdot \bm{x} = (ax_1,ax_2,\ldots,ax_n)$，$b \cdot \bm{x} = (bx_1,bx_2,\ldots,bx_n)$と定義すればよい．
  \end{proof}
\end{leftbar}

\begin{definition}{ベクトルの内積}{ベクトルの内積}
  2つのベクトル
  \[
    \bm{x} = (x_1,x_2,\ldots,x_n) \in \mathbb{R}^n,\quad \bm{y} = (y_1,y_2,\ldots,y_n) \in \mathbb{R}^n
  \]
  の内積を，
  \[
    \langle \bm{x} , \bm{y} \rangle \coloneqq x_1y_1+x_2y_2+\dotsb+x_ny_n \in \mathbb{R}
  \]
  とする．
\end{definition}

\begin{theorem}{ベクトルの内積の性質}{ベクトルの内積の性質}
  $\bm{x},\bm{y},\bm{z} \in \mathbb{R}^n$，$a \in \mathbb{R}$に対して，次の性質が成り立つ．
  \begin{enumerate}
    \item $\langle \bm{x} , \bm{y} \rangle = \langle \bm{y} , \bm{x} \rangle$
    \item $\langle \bm{x} , \bm{y} + \bm{z} \rangle = \langle \bm{x} , \bm{y} \rangle + \langle \bm{x} , \bm{z} \rangle$
    \item $\langle a \cdot \bm{x} , \bm{y} \rangle = a \cdot \langle \bm{x} , \bm{y} \rangle$
    \item $\langle \bm{x} , \bm{x} \rangle \geqq 0$，$\langle \bm{x} , \bm{x} \rangle = 0 \iff \bm{x} = \bm{0}$
    \item $\abs{\langle \bm{x} , \bm{y} \rangle} \leqq \sqrt{\langle \bm{x} , \bm{x} \rangle} \sqrt{\langle \bm{y} , \bm{y} \rangle}$
    \item $\langle \bm{x} , \bm{y} \rangle = 0 \iff \bm{x} \perp \bm{y}$
  \end{enumerate}
\end{theorem}
\phantomsection
\subsection{複素数体の定義}

\begin{definition}{$\mathbb{R}^2$の和と積}{R2の和と積}
  $(a,b),(c,d) \in \mathbb{R}^2$に関して，加法と乗法を
  \begin{align*}
    & (a,b)+(c,d)=(a+c,b+d), \\
    & (a,b) \cdot (c,d) = (ac-bd,ad+bc)
  \end{align*}
  と定める．
\end{definition}

\begin{theorem}{}{}
  \deref{R2の和と積}における$\mathbb{R}^2$は可換体である．
  乗法の単位元は$(1,0)$であり，乗法の逆元に関して$(a,b)^{-1}=(a/(a^2+b^2),-b/(a^2+b^2))$である．ただし$(a,b) \ne (0,0)$とする．
\end{theorem}

\begin{tleftbar}
\begin{proof}
  乗法の単位元に関して，
  \[
    (a,b)(1,0) = (a \cdot 1 - b \cdot 0 , a \cdot 0 + b \cdot 1) = (a,b)
  \]
  であるから，たしかに$(1,0)$は乗法の単位元である．

  また，$(a,b) \ne (0,0)$とすると，$a^2+b^2 \ne 0$であるから，$(a,b)$の乗法の逆元は
  \[
    \left( \frac{a}{a^2+b^2},\frac{-b}{a^2+b^2} \right)
  \]
  である\footnote{計算は省略する．}．これが証明すべきことであった．
\end{proof}
\end{tleftbar}

\begin{definition}{虚数単位}{虚数単位}
  \[
    i \coloneqq (0,1)
  \]
  と定め，$i$を虚数単位とよぶ．
\end{definition}

\begin{prop}{}{}
\[
  i^2 =-1
\]
\end{prop}

\begin{tleftbar}
  \begin{proof}
    \[
      i^2 = (0,1)(0,1) = (0 \cdot 0 - 1 \cdot 1 , 0 \cdot 1 + 1 \cdot 0) = (-1,0) = -1
    \]
    である．これが証明すべきことであった．
  \end{proof}
\end{tleftbar}

\begin{definition}{複素数体}{複素数体}
  $\mathbb{C} \coloneqq \mathbb{R}^2$とし，$\mathbb{C}$の元を複素数という．
\end{definition}

\begin{prop}{複素数の性質}{複素数の性質}
  $z,w \in \mathbb{C}$のとき，
  \begin{enumerate}
    \item $\overline{z \pm w} = \overline{z} \pm \overline{w}$
    \item $\overline{zw} = \overline{z} \cdot \overline{w}$
    \item $\overline{z}z = \abs{z}^2$
    \item $\abs{zw} = \abs{z} \abs{w}$
  \end{enumerate}
  となる．
\end{prop}

\begin{tleftbar}
  \begin{proof}
    1，2は演算の定義から明らか．
    3に関して，
    \[
      \overline{z}z = (a,-b)(a,b) = (a^2+b^2,0) = \abs{z}^2
    \]
    である．
    4に関して，
    \[
      \abs{zw}^2 = \overline{zw}zw = \overline{z}z \overline{w}w = \abs{z}^2 \abs{w}^2
    \]
    であるから，$\abs{zw} = \abs{z} \abs{w}$である．
  \end{proof}
\end{tleftbar}

\newpage

\phantomsection
\subsection{点列}

\begin{definition}{点列}{点列}
  写像
  \[
    a \colon \mathbb{N} \ni m \longmapsto a_m \in \mathbb{R}^n
  \]
  を$\mathbb{R}^n$の点列という．
\end{definition}

\begin{definition}{$\mathbb{R}^n$における有界性}{Rnにおける有界性}
  $\mathbb{R}^n$の部分集合$A$が有界であるとは，ある$M>0$が存在して，任意の$a \in A$に対して，
  \[
    \abs{a} \leqq M
  \]
  が成り立つことである．
\end{definition}

\begin{definition}{点列の収束}{点列の収束}
  点列$(a_m)_{m \in \mathbb{N}}$が$a \in \mathbb{R}^n$に収束するとは，任意の$\varepsilon >0$に対して，ある$N \in \mathbb{N}$が存在して，任意の$m \in \mathbb{N}$に対して，
  \[
    m \geqq N \implies \abs{a_m -a} < \varepsilon
  \]
  が成り立つことである．
\end{definition}

\begin{definition}{$\mathbb{R}^n$におけるコーシー列}{Rnにおけるコーシー列}
  点列$(a_m)_{m \in \mathbb{N}}$が$\mathbb{R}^n$においてコーシー列であるとは，任意の$\varepsilon >0$に対して，ある$N \in \mathbb{N}$が存在して，任意の$m,n \in \mathbb{N}$に対して，
  \[
    m,n \geqq N \implies \abs{a_m -a_n} < \varepsilon
  \]
  が成り立つことである．
\end{definition}

以下では$\mathbb{R}^n$の点列$(a_m)_{m \in \mathbb{N}}$を
\[
  a_m = (a_{m,1},a_{m,2},\ldots,a_{m,n})
\]
と表記することがある．

\begin{theorem}{}{}
  $\mathbb{R}^n$の点列$(a_m)_{m \in \mathbb{N}}$に関して，次のことが成り立つ．
\begin{align*}
  & \lim_{m \to \infty} a_m = a \in \mathbb{R}^n \\
  \iff & \lim_{m \to \infty} a_{m,i} = a_i \in \mathbb{R} \quad (i=1,2,\ldots,n)
\end{align*}
\end{theorem}

\begin{leftbar}
\begin{proof}
  必要条件であること，十分条件であることを分けて証明する．

  必要条件であることを示す．$\lim_{m \to \infty} a_m = a \in \mathbb{R}^n$とする．
  このとき，任意の$\varepsilon >0$に対して，ある$N \in \mathbb{N}$が存在して，任意の$m \in \mathbb{N}$に対して，
  \[
    m \geqq N \implies \abs{a_m -a} < \varepsilon
  \]
  が成り立つ．
  ここで，$m \geqq N$とすると，
  \[
    \abs{a_m -a} = \left \{ \sum_{i=1}^{n} (a_{m,i} -a_i)^2 \right \}^\frac{1}{2}  < \varepsilon
  \]
  であるから，$\abs{a_{m,i} -a_i} < \varepsilon$である．
  よって，$\lim_{m \to \infty} a_{m,i} = a_i \in \mathbb{R} \quad (i=1,2,\ldots,n)$である．

  十分条件であることを示す．$\lim_{m \to \infty} a_{m,i} = a_i \in \mathbb{R} \quad (i=1,2,\ldots,n)$とする．
  任意の$\varepsilon >0$に対して，ある$N_i \in \mathbb{N}$が存在して，任意の$m \in \mathbb{N}$に対して，
  \[
    m \geqq N_i \implies \abs{a_{m,i} -a_i} < \frac{\varepsilon}{\sqrt{n}}
  \]
  が成り立つ．
  ここで，$N \coloneqq \max \{ N_1,N_2,\ldots,N_n \}$とすると，$m \geqq N$のとき，
  \[
    \abs{a_m -a} =\left \{ \sum_{i=1}^{n} (a_{m,i} -a_i)^2 \right \}^{\frac{1}{2}}  < \left \{ \sum_{i=1}^{n} \left( \frac{\varepsilon}{\sqrt{n}} \right)^2 \right \}^\frac{1}{2} = \varepsilon
    \]
    となり，$\lim_{m \to \infty} a_m = a \in \mathbb{R}^n$である．これが証明すべきことであった．
\end{proof}
\end{leftbar}

\begin{theorem}{}{}
  $\mathbb{R}^n$の点列$(a_m)_{m \in \mathbb{N}}$に関して，次のことが成り立つ．
\begin{align*}
  & (a_m)_{m \in \mathbb{N}} ~{:}  \text{コーシー列} \\
  \iff & (a_{m,i})_{m \in \mathbb{N}} ~{:} \text{コーシー列}\quad (i=1,2,\ldots,n)
\end{align*}
\end{theorem}

\begin{leftbar}
  \begin{proof}
    必要条件であること，十分条件であることを分けて証明する．

    必要条件であることを示す．$(a_m)_{m \in \mathbb{N}}$がコーシー列であるとする．
    このとき，任意の$\varepsilon >0$に対して，$N \in \mathbb{N}$が存在し，任意の$m,l \in \mathbb{N}$に対して
    \[
      m,l  \geqq N \implies \abs{a_m - a_l}<\varepsilon
    \]
    が成り立つ．

    ここで，
    \[
      \abs{a_m - a_l} = \left \{ \sum_{i=1}^{n} \abs{a_{m,i}-a_{l,i}}^2 \right \}^\frac{1}{2} < \varepsilon
    \]
    であるから，$ l,m \geqq N$のとき$\abs{a_{m,i}-a_{l,i}}<\varepsilon$~$(i=1,2,\ldots,n)$であり，ここまでで必要条件であることが示された．

    十分条件であることを示す．$(a_{m,i})_{m \in \mathbb{N}}$~$(i=1,2,\ldots,n)$がコーシー列であるとする．

    このとき，任意の$\varepsilon /\sqrt{n}>0$に対して，ある$N_i \in \mathbb{N}$が存在して，任意の$m,l \in \mathbb{N}$に対して，
    \[
      m ,l\geqq N_i \implies \abs{a_{m,i}-a_{l,i}}<\frac{\varepsilon}{\sqrt{n}}
    \]
    である．ここで，$N \coloneqq \max \{ N_1 ,N_2,\ldots,N_n \}$とすると，$m,l \geqq N$のとき
    \[
      \abs{a_m - a_l }=\left  \{ \sum_{i=1}^{n} \abs{a_{m,i}-a_{l,i}}^2 \right \}^\frac{1}{2} < \left \{ \sum_{i=1}^{n} \left (  \frac{\varepsilon}{\sqrt{n}} \right )^2 \right \}^\frac{1}{2}=\varepsilon
    \]
    となり，ここまでで十分条件であることも示された．これが証明すべきことであった．
  \end{proof}
\end{leftbar}

\begin{theorem}{}{}
  $\mathbb{R}^n$の点列$(a_m)_{m \in \mathbb{N}}$に関して，次のことが成り立つ．
\begin{align*}
  & (a_m)_{m \in \mathbb{N}}~ {:}  \text{収束} \\
  \iff & (a_m)_{m \in \mathbb{N}} ~{:}  \text{コーシー列}
\end{align*}
\end{theorem}

\begin{leftbar}
  \begin{proof}
    これまでに示した事実を並べると，
\begin{align*}
  & (a_m)_{m \in \mathbb{N}}~ {:} \text{収束} \\
  \iff & (a_{m,i})_{m \in \mathbb{N}}~ {:} \text{収束} \quad (i=1,2,\ldots,n) \\
  \iff & (a_{m,i})_{m \in \mathbb{N}}~{:} \text{コーシー列}\quad (i=1,2,\ldots,n)  \\
  \iff & (a_m)_{m \in \mathbb{N}} ~{:} \text{コーシー列}
\end{align*}
となる．これが証明すべきことであった．
\end{proof}
\end{leftbar}

\begin{theorem}{}{}
  $\mathbb{R}^n$の点列$(a_m)_{m \in \mathbb{N}}$とその部分列$(a_{m(k)})_{k \in \mathbb{N}}$に関して，次のことが成り立つ．
  \begin{align*}
    &\lim_{m \to \infty} a_m = b \in \mathbb{R}^n \\
    \implies & \lim_{k \to \infty} a_{m(k)} = b \in \mathbb{R}^n
  \end{align*}
\end{theorem}

\begin{leftbar}
\begin{proof}
  任意の$\varepsilon >0$に対して，ある$N \in \mathbb{N}$が存在して，任意の$m \in \mathbb{N}$に対して，
  \[
    m \geqq N \implies \abs{a_m -b} < \varepsilon
  \]
  が成り立つ．
  ここで，$k \in \mathbb{N}$に対して，$m(k) \geqq k$であるから，
  \[
    k \geqq N \implies m(k) \geqq N \implies \abs{a_{m(k)} -b} < \varepsilon
  \]
  であるから，$\lim_{k \to \infty} a_{m(k)} = b \in \mathbb{R}^n$である．これが証明すべきことであった．
\end{proof}
\end{leftbar}

\begin{theorem}{}{}
  $\mathbb{R}^n$の点列$(a_m)_{m \in \mathbb{N}}$に関して，次のことが成り立つ．
  \begin{align*}
    & (a_m)_{m \in \mathbb{N}}~ {:}  \text{収束} \\
    \implies  & (a_m)_{m \in \mathbb{N}} ~{:}  \text{有界}
  \end{align*}
\end{theorem}

\begin{leftbar}
  \begin{proof}
    各$k \in \{ 1,2,\ldots,n \}$に対し$(a_{m,k})_{m \in \mathbb{N}}$は収束するので，
    先の命題により，$(a_{m,k})_{m \in \mathbb{N}}$は有界である．
    よって，$(a_m)_{m \in \mathbb{N}}$も有界である．これが証明すべきことであった．
  \end{proof}
\end{leftbar}
